% Artículo VII — primera redacción editorial, 10 de septiembre de 2026.
% Apertura integrada con las demostraciones de la revisión 06.
% Procedencia y condiciones: PLAN_RESIDENCIAS_DEMOSTRATIVAS.md.
% Las referencias internas remiten a demostraciones incluidas en el cuerpo.

\begin{abstract}
Se desarrolla la relación entre memoria orientada, geometría gravitatoria
y evolución cosmológica desde la construcción APP--TRIT--TPK. El recorrido
demostrativo articula la ley de dilución \(a^{-6}\), la existencia y
unicidad del umbral de rebote en la realización homogénea especificada y la
persistencia local del rebote transversal bajo perturbaciones controladas.
La memoria conserva orientación, acarreo e historia durante el retorno de
fase; su transporte admite realizaciones operatorias y geométricas que
preservan esas distinciones. El paso hacia la dinámica gravitatoria reúne
la corriente de pantalla, el funcional de amplitud torsional, su
normalización y su compatibilidad con el refinamiento. En la acción mínima
de Cartan--Holst, la torsión queda determinada algebraicamente por la
corriente de espín y evoluciona con la geometría y los campos materiales.
Los balances cosmológicos conservan las condiciones de las fuentes y de
los acoplamientos en cada carta. El desarrollo se completa con la
información de horizonte, las operaciones de lectura del observador y la
respuesta relacional del estado conjunto. Se obtiene explícitamente el
funcional de amplitud desde la acción espinorial, la eliminación de la
contorsión y la variación métrica, con un dominio material de signo
positivo y conservación demostrados. Para la realización plana de polvo
se integra la evolución en todo tiempo y se prueba la completitud de sus
geodésicas causales. La propagación bidireccional se construye mediante
inversas avanzadas y retardadas del transporte de rutas y una condición
de compatibilidad en la frontera.
\end{abstract}
\clearpage

\section{Introducción}
\label{vii:sec:introduccion}

La relación entre gravitación y memoria adquiere contenido matemático
cuando se especifican el estado que conserva la historia, la operación que
la transporta y la ecuación en la que interviene su realización geométrica.
En la construcción considerada, el retorno de fase conserva una
periodicidad observable mientras el registro incorpora orientación,
acarreo y profundidad. El problema gravitatorio se sitúa en el transporte
de esa información y en su contribución a las variables geométricas y
materiales de la evolución.

Tres resultados organizan la exposición: la dilución de la memoria, el
rebote de la realización homogénea y la persistencia local de su umbral
transversal. La primera relación se expresa mediante dos lectores de una
misma ruta,
\begin{equation}
 \frac{a_n}{a_{\mathrm{ref}}}=\varphi^{n/3},
 \qquad M_n=\varphi^{-2n},
 \qquad
 M_n=\left(\frac{a_n}{a_{\mathrm{ref}}}\right)^{-6}.
 \label{vii:eq:programa-dilucion}
\end{equation}
Aquí \(a_{\mathrm{ref}}>0\) fija la normalización del factor de escala.
La potencia expresa la relación entre el crecimiento de la realización
espacial y el carácter cuadrático del lector de memoria. La contribución
gravitatoria requiere, además, determinar la corriente y la amplitud que
acompañan a esa dependencia de escala. Ambas construcciones ocupan una
sección propia antes de introducirlas en las ecuaciones cosmológicas.

La Holografía Modular Triádica aporta el soporte y las operaciones de esta
genealogía. APP conserva las hojas aditiva y multiplicativa y sus datos
residuales; TRIT tipa los regímenes y las orientaciones; TPK transporta,
memoriza y prolonga el estado. Sus publicaciones internas transmiten a
las construcciones posteriores tanto los valores como las relaciones que
los determinan. La Mecánica Dimensional realiza esas estructuras mediante
secciones de acción, magnitudes geométricas y observables físicos.

El primer movimiento del artículo reúne construcción y geometría. La
descomposición del transporte distingue memoria simétrica y circulación
orientada. La representación del grupoide de rutas en geometría de Cartan
se desarrolla con composición, inversión, covariancia y refinamiento. La
acción, el parámetro de Barbero--Immirzi, la constante gravitatoria y la
unidad areal conservan sus antecedentes estructurales; su reutilización
transmite las condiciones que individualizan cada sección.

El segundo movimiento desarrolla amplitud, evolución y persistencia.
La corriente de pantalla se lleva a un diferencial respecto de la
conexión mediante la incidencia transportada y la extensión minimizante.
La realización espinorial se desarrolla, a su vez, desde las matrices
internas de Clifford y su acción material. Ambas composiciones conservan
el funcional que las define y sus respectivas reglas de variación.
La normalización tracial
fija la lectura de densidad correspondiente y se mantiene diferenciada de
una suma extensiva. La variación de la conexión determina la respuesta
torsional de la acción mínima. Al evolucionar la cotetrada y los campos
materiales, cambia también la torsión algebraicamente asociada a su
corriente de espín.

La amplitud del sector espinorial queda determinada por
\begin{equation}
 s_0[\varrho,V_c]
 =\frac{3\kappa c^2\hbar^2}{16}
   \frac{\gamma^2}{1+\gamma^2}
   \frac{\operatorname{Tr}
    (\varrho\widehat{\mathscr Q}_{\rm sp})}{V_c^2},
 \qquad \kappa=\frac{8\pi G}{c^4}.
 \label{vii:eq:programa-amplitud}
\end{equation}
Aquí \(\varrho\) es el estado material normalizado, \(V_c\) el volumen
comóvil y \(\widehat{\mathscr Q}_{\rm sp}\) el observable cuártico
construido a partir de la corriente. El
teorema~\ref{vii:thm:amplitud-material-explicita} deriva la fórmula y
el corolario~\ref{vii:cor:dominio-positivo-conservado} exhibe estados
en los que la amplitud es positiva y permanece constante. La dependencia
del segundo momento conserva información del estado incluso cuando la
media de la corriente se anula. El dato material y su volumen fijan una
realización; las constantes de acoplamiento proceden de las secciones
estructurales construidas previamente.

La eliminación de la torsión se trata también para la acción de
extensión mínima, cuya dependencia respecto de la conexión es, en
general, no lineal. Su diferencial determina la ecuación estacionaria;
el Hessiano controla la continuación local y una identidad integral
determina la acción efectiva. Esta formulación permite conservar la
dependencia completa del mínimo al efectuar las variaciones.

La formulación del rebote conserva las densidades positivas, el signo de
la contribución torsional y la desigualdad de exponentes de la realización
homogénea. Su persistencia local se demuestra mediante dos controles
separados: una cota uniforme de la perturbación conserva el cambio de
signo en los extremos de un intervalo; una cota uniforme de su derivada
conserva la monotonía y la transversalidad. Las realizaciones de los
selectores cosmológicos se enuncian con sus propios dominios, junto a los
invariantes y balances que producen.

La solución plana de polvo proporciona además un desarrollo global:
\[
 a(t)=a_{\rm b}\left(1+
       \frac{(t-t_{\rm b})^2}{\tau^2}\right)^{1/3},
 \qquad
 a_{\rm b}^3=\frac{s_0}{\rho_0},\qquad
 \tau^2=\frac{s_0}{6\pi G\rho_0^2},
 \quad c=1.
\]
La prueba establece regularidad de la curvatura y completitud causal
en esa realización homogénea. Su integración y la persistencia local
del umbral describen dos propiedades complementarias, con sus
respectivos dominios explícitos.

La conservación del tensor residual se formula en la carta métrica
seleccionada, con acoplamientos constantes sobre ella y una fuente
material covariantemente conservada. Su descomposición en traza y parte
sin traza permite identificar el intercambio entre ambos sectores;
la traza constante caracteriza el caso de conservación separada.

El tercer movimiento estudia horizontes y respuesta relacional. El estado
reducido, los registros accesibles y la purificación permiten describir
qué información publica una lectura y qué información permanece en las
correlaciones del estado conjunto. La definición operacional del
observador comprende orientación y trivialización del transporte,
restricción a observables y proyección al subsistema. La respuesta
autoinercial se conecta con esas operaciones mediante la factorización
por la traza global y el defecto geométrico de proyección.

Los balances de horizonte conservan igualmente las convenciones de
área, energía, temperatura y estado que utilizan.
La propagación avanzada y retardada se compone sobre el transporte
unitario de rutas. La igualdad de sus lecturas en la frontera se
expresa mediante una restricción lineal de fuentes; su proyección
determina la respuesta variacional en ese dominio.

Cada desarrollo culmina con el objeto construido, la operación que lo
determina, su invariante y su función física dentro de la realización.
El orden de las pruebas sigue las dependencias de esos objetos; la
organización narrativa hace visible su relación con los tres resultados
principales desde el comienzo. Las conclusiones reúnen también la
operación del observador y la respuesta relacional, de modo que los
balances informacionales y la dinámica gravitatoria conserven su
conexión explícita.
